\subsection{Truth Tables}

\content{
\begin{definition} We use the following symbols for conjunction (and),
disjunction (or), and negation (not): 
$$ \text{and: } \land,\quad \text{or: }\lor,\quad 
\text{not: }\neg. $$
Sometimes you may see the symbol $\sim$ used for negation instead of $\neg$.
\end{definition}
}{% presentation
}{% nopresentation
}{% comments
}

\content{
\begin{example} Let $C$ be the statement ``I like Coca-Cola" and $P$ be the
statement ``I like Pepsi". Write the following in symbols.
\begin{enumerate}
\item I like Coke and Pepsi
\item I like Coke or I like Pepsi
\item I do not like Coke
\item I like Coke and I do not like Pepsi
\end{enumerate}
\end{example}
}{% presentation
\vspace{1in}
}{% nopresentation
}{% comments
}

\content{
\begin{example} Let $R$ be the statement ``This couch is made of leather", and
$B$ be the statement ``This couch reclines". Write the following in symbols.
\begin{enumerate}
\item This couch is made of leather, but it does not recline
\item This couch is not made of leather, but it relines
\item This couch does not recline, and is not made of leather
\end{enumerate}
\end{example}
}{% presentation
\vspace{1in}
}{% nopresentation
}{% comments
}

\content{
\begin{definition} A \textbf{truth table} is a table showing the value of a
compound statement for all possible values of the simple statements. 
\end{definition}
}{% presentation
}{% nopresentation
}{% comments
}

\content{
\begin{example} Construct a truth table for the compound statement $p\land q$.\\
\begin{tabular}{|c|c|c|}
\hline
$p$ & $q$ &\hspace{0.75cm}  \\
\hline
T & T &  \\
\hline
T & F &  \\
\hline
F & T &  \\
\hline
F & F &  \\
\hline
\end{tabular}
\end{example}
}{% presentation
}{% nopresentation
}{% comments
}

\content{
\begin{example} Your turn! Construct a truth table for the compound statement
$p\lor q$.\\
\begin{tabular}{|c|c|c|}
\hline
$p$ & $q$ &\hspace{0.75cm}  \\
\hline
T & T &  \\
\hline
T & F &  \\
\hline
F & T &  \\
\hline
F & F &  \\
\hline
\end{tabular}
\end{example}
}{% presentation
}{% nopresentation
}{% comments
}

\content{
\begin{example} Construct a truth table for the compound statement $p\lor \neg
q$.\\
\begin{tabular}{|c|c|c|c|}
\hline
$p$ & $q$ &\hspace{0.75cm} & \hspace{0.75cm} \\
\hline
T & T & & \\
\hline
T & F & & \\
\hline
F & T & & \\
\hline
F & F & & \\
\hline
\end{tabular}
\end{example}
}{% presentation
}{% nopresentation
}{% comments
}

\content{
\begin{example} Construct a truth table for the compound statement 
$\neg (p\land q)$.\\
\begin{tabular}{|c|c|c|c|}
\hline
$p$ & $q$ &\hspace{0.75cm} & \hspace{0.75cm} \\
\hline
T & T & & \\
\hline
T & F & & \\
\hline
F & T & & \\
\hline
F & F & & \\
\hline
\end{tabular}
\end{example}
}{% presentation
}{% nopresentation
}{% comments
}

\content{
\begin{example} Construct a truth table for the compound statement 
$p\land \neg(q\lor r)$.\\
\begin{tabular}{|c|c|c|c|c|c|}
\hline
$p$ & $q$ & $r$ & \hspace{0.75cm} & \hspace{0.75cm} & \hspace{0.75cm} \\
\hline
T & T & T & & & \\
\hline
T & T & F & & &\\
\hline
T & F & T& & & \\
\hline
T & F & F& & &\\
\hline
F & T & T & & &\\
\hline
F & T & F& & & \\
\hline
F & F & T& & &\\
\hline
F & F & F& & &\\
\hline
\end{tabular}
\end{example}
}{% presentation
}{% nopresentation
}{% comments
}

\content{
\begin{definition}
A conditional is a logical compound statement in which a statement $p$, implies
a statement $q$. This is written $p\Rightarrow q$ (or sometimes $p\rightarrow
q$). The truth table of the conditional $p\Rightarrow q$ is:\\
\begin{tabular}{|c|c|c|}
\hline
$p$ & $q$ & $p \Rightarrow q$ \\
\hline
T & T &  T\\
\hline
T & F &  F\\
\hline
F & T &  T\\
\hline
F & F &  T\\
\hline
\end{tabular}
\end{definition}
}{% presentation
}{% nopresentation
}{% comments
}

\content{

}{% presentation
}{% nopresentation
}{% comments
}
